% Planned analysis, resource trade-offs, and quality assurance.
\subsection{Data Analysis and Quality Assurance}

Time to full coverage will be right-censored at the common timeout. The main figure, to be produced after the experiments, will plot restricted mean time to completion up to that limit against target mean fade duration, separately for balanced and unequal progress. It will include the fixed partition, reactive coordination, non-predictive long backlog, and the two predictive variants. The primary contrast is trend-based prediction (S4-T) minus reactive coordination (S2), with 95\% paired bootstrap confidence intervals obtained by resampling matched scenario seeds. S4-T versus S3 and S4-T versus S4-P are pre-specified secondary comparisons, not alternative primary tests chosen after inspecting results.

Realized outage duration and link availability will be reported by configuration alongside completion proportion and coverage at timeout. Secondary measures will include redundant observations, travel distance, idle time, and pre-outage reservations that were successfully committed. These distinguish useful coordination from warnings that did not change an assignment or could not be acted upon. Distance-related stalls and budget failures will be reported separately from radio fades. If performed, the optional distance-loss sensitivity check will assess whether conclusions depend on route-dependent connectivity.

Completion is the first time every target cell has at least one observation, excluding return and data upload. A redundant observation is a new cell-entry observation of a previously observed cell, not every simulator tick spent inside it. Counts, distance, and idle time will be accumulated until completion or timeout and interpreted alongside achieved coverage; a stalled run must not appear efficient merely because it travels or communicates less. With only administrative censoring at the common timeout, the restricted-mean estimate is
\begin{equation}
\label{eq:restricted-mean}
\widehat{\mu}_{\tau}=\frac{1}{N}\sum_{s=1}^{N}\min(T_s,\tau).
\end{equation}
Here $N$ is the number of evaluated seeds for a configuration within one condition, $s$ indexes those seeds, $T_s$ is full-coverage time in seconds, and $\tau=7200$ s is the common stopping limit. For an unfinished mission, treat $T_s$ as beyond $\tau$: its contribution is $\tau$, not a claimed completion. Resample whole paired seeds 2000 times within each experimental condition for 95\% bootstrap intervals. Report differences in seconds, relative changes against the reference, and completion proportions rather than relying only on interval inclusion of zero. Software crashes are implementation failures to investigate and rerun, not mission timeouts to insert into this estimate.

For warning evaluation, a timely alert occurs while connected one to three exchange intervals before a radio-outage onset. Match alerts chronologically to the first eligible subsequent onset, at most once per event; collapse consecutive positive forecasts into one alert until the warning becomes negative or an outage/reconnection completes the episode. Unmatched alerts are false alarms and unmatched eligible onsets are misses. Exclude and count boundary cases lacking a complete observation window: the first three opportunities for onset recall and the last four for alert matching. Report event counts, precision, recall, F1, and lead-time distributions. Precision with no alerts and recall with no eligible outages are undefined; report these cases explicitly. F1 is zero when eligible outages exist but there are no matched alerts, and undefined when there are no eligible outages. Fixed-trace replay isolates warning quality; paired mission runs measure consequences under route-dependent connectivity. Better warning scores alone do not establish better mission performance.

Count attempted and delivered messages, serialized bytes, and budget-exhausted sessions for every configuration, including failed attempts. Use the same documented message encoding so the byte counts can be reproduced. Measure CPU time separately for forecasting, allocation, and resynchronization, and report both total coordination CPU time per run and time per decision. Repeated timings on saved inputs will distinguish small computation costs from timer noise. Record hardware, operating system, Python/package versions, repetitions, and the timing procedure; exclude plotting and file-writing overhead from component timings.

To compare algorithmic workload without choosing a particular USV processor, also count arithmetic operations separately for forecasting, allocation, and resynchronization. Document counts from the implemented pseudocode and executed branches, and check them against small hand-worked cases in the pilot. Count each floating-point addition, subtraction, multiplication, or division as one operation; a fused multiply-add, if used, counts as two. Report comparisons and logarithm/square-root calls separately because treating all operation types as equally costly would be misleading. State numerical precision and any library-kernel counting assumptions. These are algorithmic counts, not measurements of Python interpreter or processor instructions. For component $c$, report both the raw count $O_c$ accumulated until completion or timeout and its billion-operation equivalent,
\begin{equation}
\label{eq:gop-count}
G_c=10^{-9}O_c.
\end{equation}
Here $c$ identifies forecasting, allocation, or resynchronization, and $G_c$ is a workload in GOPs (billions of counted arithmetic operations), not GOP/s throughput. Small workloads will retain raw counts rather than be rounded to zero. Report per-decision counts and totals alongside achieved coverage; exclude simulator motion, plotting, logging, and profiling overhead from coordination-operation counts.

Measure peak resident process memory (RAM) in isolated runs using the same host, software environment, and bounded logging buffers. Report this as total simulator-process memory, not memory attributable solely to a predictor. Separate profiling runs will measure incremental Python allocations for forecasting and allocation, with the interpreter and shared state as the baseline; these measurements exclude some native-library allocations and will be labelled accordingly. Memory profiling and CPU benchmarking will be performed separately to avoid profiling overhead distorting the timing comparison.

The evaluation criteria are mission effectiveness, warning quality, and resource cost. Lower restricted mean completion time indicates a time benefit, interpreted together with completion probability and uncertainty. We will compare time savings with additional transmitted bytes, arithmetic-operation counts/GOPs, coordination CPU time, and peak RAM in tables and trade-off plots. If a configuration is faster but more costly, both effects will be reported; no universal winner or arbitrary weighted score will be imposed. The comparison with S3 will indicate whether the warning adds value beyond a longer queue alone.

Map accuracy requires a separate qualification. The model assumes correct observations of visited cells and contains no depth estimates, sonar noise, or positioning error. Full coverage therefore does not establish accurate bathymetry. As a coordination diagnostic, let $C(t)$ be the set of cells observed by either USV at simulation time $t$ seconds and $C_i(t)$ the coverage known to USV $i\in\{1,2\}$. Define the missing-information fraction as
\begin{equation}
\label{eq:missing-coverage}
D_i(t)=\frac{|C(t)\setminus C_i(t)|}{M}.
\end{equation}
Here $M=240$ is the number of target cells, $\setminus$ denotes set difference, and $|\cdot|$ denotes the number of cells in a set. Record $D_i(t)$ at scheduled exchanges before synchronization and report its average over those observations during each active run and its value at stopping. This measures missing local coverage information, not seabed reconstruction error. Under ideal sensing, locally known coverage outside $C(t)$ is an implementation error. Physical map accuracy would require sensor and localization models or measured survey data, which are outside this study.

Resource trade-offs will be assessed by comparing mission benefits with communication and computation costs. Gong et al. motivate resource-aware coordination but do not supply an interchangeable power model for our vessels \cite{gong2022resource}. The core analysis will use arithmetic-operation counts/GOPs, message attempts and bytes, host CPU time, and memory as resource indicators, without requiring the power model of a particular USV. Operation counts and communication volume support discussion of workload trends; host timings and memory measurements describe this implementation. They do not establish a universal energy ordering: operation type, numerical precision, memory traffic, processor/radio architecture, duty cycle, and propulsion can change the energy cost. A shorter mission may reduce some operating costs while prediction adds communication or computation, so both sides of that trade-off will be reported. We will not convert these indicators into joules, battery endurance, battery ageing, or device lifetime. Actual vessel-energy savings require separate platform measurements and remain outside the measured outcomes of this study.

Quality assurance includes deterministic connected/forced-outage/reconnection traces, unique-owner and message-budget assertions, tests for no future-information leakage, reproducible seeds, and reconstruction of metrics from raw logs. Both authors reproduce a sample run and review failed/unfinished cases. Report task conflict counts and redundant observations even if they are zero; do not create a conflict-prone benchmark to manufacture an effect. These checks verify implementation, not environmental fidelity.
